\documentclass{article}

\usepackage[preprint]{neurips_2025}

\usepackage{amsmath,amssymb}
\usepackage{booktabs}
\usepackage{graphicx}
\usepackage{hyperref}

\title{Post-Training Typed-Decision Models with Proper-Scoring Rewards:\\Accuracy, Calibration, and the Underconfidence Effect}
\author{Sai Krishna Rallabandi}

\begin{document}
\maketitle

\begin{abstract}
Calibrated decisions require probabilities that mean what they say, yet the
post-training recipes that maximise decision accuracy routinely damage
calibration, and those that preserve it leave accuracy on the table. We show
both effects can be had in sequence: post-training a 4B prefill-only
decision model with soft-target cross-entropy plus Brier and
ranked-probability-score terms raises held-out decision accuracy 0.665 to
0.800 and cuts KL from gold threefold, while expected calibration error
regresses 0.091 to 0.162. A reliability analysis traces the regression to
systematic \emph{underconfidence} induced by distributional supervision, and
a single temperature (T*~=~0.35, fit on disjoint training data) restores ECE
to 0.031 with accuracy unchanged. The result is a small, single-pass model
that beats an accuracy-first specialist and a calibration-first baseline on
both axes at once.
\end{abstract}

\section{Introduction}

Autonomous systems increasingly act on model outputs without a human in the
loop: routing tickets, triaging incidents, approving invoices. Each such act
is a \emph{typed decision} — a choice among named options, a score on an
ordered scale, or a yes/no verdict — reported as a probability distribution.
The distribution is the product: downstream thresholds (automate, defer,
escalate) are only meaningful if 0.80 means right 80\% of the time. The
problem is important because miscalibrated automation either acts when it
should defer or defers when it could act, and both failures compound at
machine scale.

The problem is hard because the standard post-training objectives pull in
opposite directions. Preference-based RL (RLHF) rewards text that
\emph{sounds} right and narrows output distributions, the wrong objective
for unattended decisions. Verifiable-reward RL (RLVR) maximises accuracy
with binary rewards and drives models toward extreme overconfidence; across
decision benchmarks it wins accuracy while wrecking calibration. The naive
fix — supervised fine-tuning — calibrates better but gains less accuracy.
No single-stage recipe in common use delivers both.

It had not been solved in this setting because the two fixes live in
different stages: proper-scoring rewards shape \emph{distributions} during
training, while temperature scaling reshapes \emph{confidence} after it, and
prior work evaluates one or the other rather than composing them on a fixed
decision benchmark with distributional supervision. Our approach is the
composition: (i)~post-train a prefill-only decision model with
cross-entropy against soft teacher-averaged targets plus Brier and RPS
calibration terms, then (ii)~fit a single temperature on disjoint data.
On 400 held-out business-decision cases the composed model reaches 0.800
decision accuracy with ECE 0.031, against 0.665/0.091 for its base, 0.767
with ECE 0.214 for an accuracy-first specialist, and 0.532 with ECE 0.085
for a calibration-first baseline.

\begin{itemize}
\item We show proper-loss post-training on soft decision targets improves
accuracy, proper scores, and selective prediction together (Section~5).
\item We identify an \emph{underconfidence effect}: distributional
supervision spreads mass and depresses argmax confidence below argmax
accuracy, which ECE punishes symmetrically (Section~6).
\item We show a single honestly-fit temperature resolves the effect with
accuracy unchanged, yielding a model that dominates both the
accuracy-first and calibration-first baselines on both axes (Section~5).
\end{itemize}

\begin{figure}[t]
\centering
\begin{minipage}{0.92\linewidth}\small
\textbf{Running example (invoice\_processing\_000008, held out).} Invoice
INV-2026-6038 (\$36,720, Acme Fabrication) against PO-4946; delivery
accepted with exceptions, received quantity 0. Gold vs.\ model:
\texttt{discrepancy\_severity} (score): gold $\{3{:}0.83\}$ vs.\ model
$\{3{:}0.62, 2{:}0.30\}$; \texttt{disposition} (choice): gold
$\{\text{manual\_review}{:}0.57, \text{hold}{:}0.34\}$ vs.\ model
$\{\text{manual\_review}{:}0.50, \text{hold}{:}0.43\}$;
\texttt{duplicate} (noul): gold $\{\text{false}{:}0.96\}$ vs.\ model
$\{\text{false}{:}0.95\}$; \texttt{matches\_order} (noul): gold
$\{\text{false}{:}0.92\}$ vs.\ model $\{\text{false}{:}0.89\}$;
\texttt{urgency} (score): gold $\{2{:}0.50, 1{:}0.37\}$ vs.\ model
$\{2{:}0.56, 1{:}0.27\}$. All five argmax decisions correct; predicted mass
tracks gold mass with slightly wider spread — the signature of
Section~6's underconfidence effect.
\end{minipage}
\caption{A real held-out case used throughout: state, gold distributions,
and measured model outputs. Panels in later sections refer to its
\texttt{urgency} question (Q5) and \texttt{disposition} question (Q2).}
\label{fig:running}
\end{figure}

\section{Background / Related Work}

Two methodological commitments separate this work from its neighbours, and
everything else is credit. First, \emph{confidence is the output
distribution, not a second channel.} The RLCR recipe~\cite{damani2025rlcr}
proves that correctness plus a bounded proper rule incentivises both
accuracy and calibration, but instantiates confidence as verbalised text
produced alongside the answer; we instantiate it as the decision
distribution itself, which is what a single-pass model returns and what
downstream thresholds consume. We follow RLCR's core theoretical
constraint — Brier, not log loss, as the calibration term — while
discarding the verbalisation machinery. Second, \emph{supervision is
distributional.} Where RLVR-style work rewards discrete hits and
SFT-style work distills hard labels, we train cross-entropy against
teacher-averaged soft targets, so the Brier and RPS terms refine an
already-distributional objective rather than fighting a discretised one.
The calibration-aware RL diagnosis of~\cite{yaldiz2026calibration} —
decision tokens inherit overconfidence with no calibrated paths to
reinforce — motivates our reasoning-free, single-pass architecture, in
which there are no reasoning traces to inherit from. TypeSafe's
RLCD program~\cite{typesafe2026rlcd} states the industrial version of our
objective; it publishes no recipe, curves, or training details, so it is a
goal statement we operationalise rather than prior art we build on. We
build directly on the open Kev post-training stack (causal LM, LoRA,
pointer readout) and evaluate inside a fixed typed-decision harness.
